\subsection{\texorpdfstring{系数 \(c_{\alpha\beta\gamma}\)}{系数 c\_\{\textbackslash alpha\textbackslash beta\textbackslash gamma\}}}

令 \(\Omega\) 为所有 \((x,y)\in\mathbf{U}\times\mathbf{U}\) 的集合，其中 \(x.y\) 有定义且属于 \(\mathbf{U}\)。则 \(\Omega\) 是 \(\mathbf{U}\times\mathbf{U}\) 中的开集，且映射 \((x,y)\mapsto x.y\) 从 \(\Omega\) 到 \(\mathbf{U}\) 是解析的。因此，\(z_{1},\ldots,z_{n}\)（即 \(z=x.y\) 的坐标）可在原点附近按 \(x_{1},\ldots,x_{n},y_{1},\ldots,y_{n}\) 的幂展开为幂级数。因此，存在定义良好的常数 \(c_{\alpha_{1},\ldots,\alpha_{n},\beta_{1},\ldots,\beta_{n}}\in\mathbf{K}\)，使得

\[
z _ {1} ^ {\gamma_ {1}} \dots z _ {n} ^ {\gamma_ {n}} = \sum_ {\alpha_ {1},\dots,\alpha_ {n},\beta_ {1},\dots,\beta_ {n} \in \mathbf {N}} c _ {\alpha_ {1},\dots,\alpha_ {n},\beta_ {1},\dots,\beta_ {n},\gamma_ {1},\dots,\gamma_ {n}} x _ {1} ^ {\alpha_ {1}} \dots x _ {n} ^ {\alpha_ {n}} y _ {1} ^ {\beta_ {1}} \dots y _ {n} ^ {\beta_ {n}}\tag{1}
\]

对于 \(\gamma_{1},\ldots,\gamma_{n}\) 属于 \(\mathbf{N}\)。采用 Differentiable and Analytic Manifolds, R 中的约定，我们把这些公式简写为：

\[
(x. y) ^ {\gamma} = \sum_ {\alpha , \beta \in \mathbf {N} ^ {n}} c _ {\alpha , \beta , \gamma} x ^ {\alpha} y ^ {\beta} \quad (\gamma \in \mathbf {N} ^ {n}).\tag{2}
\]

由于 \(x.0 = 0.x = x\) 对 \(x\in \mathbf{U}\) 成立，

\[
c _ {\alpha , 0, \gamma} = c _ {0, \alpha , \gamma} = \delta_ {\alpha \gamma}\tag{3}
\]

其中 \(\delta_{\alpha \gamma}\) 是克罗内克指标。特别地，以下从现在起用 \(k\) 代替 \(\varepsilon_k\)，其中 \(k = 1, \ldots, n\)，

\[
(x. y) _ {k} = x _ {k} + y _ {k} + \sum_ {| \alpha | \geqslant 1, | \beta | \geqslant 1} c _ {\alpha , \beta , k} x ^ {\alpha} y ^ {\beta}.\tag{4}
\]

记 \(c_{\alpha \beta} = (c_{\alpha \beta 1}, c_{\alpha \beta 2}, \ldots, c_{\alpha \beta n}) \in \mathbf{K}^n\)，于是得到

\[
x. y = x + y + \sum_ {| \alpha | \geqslant 1, | \beta | \geqslant 1} c _ {\alpha \beta} x ^ {\alpha} y ^ {\beta}.\tag{5}
\]

在 (5) 的右端，我们取 2 次齐次分量：

\[
\mathrm{B} (x, y) = \sum_ {i, j = 1} ^ {n} c _ {i j} x _ {i} y _ {j}
\]

于是 \((x,y)\mapsto \mathrm{B}(x,y)\) 是从 \(\mathbf{K}^n\times \mathbf{K}^n\) 到 \(\mathbf{K}^n\) 的双线性映射。于是

\[
x. y \equiv x + y + \mathrm{B} (x, y) \mod \deg 3.\tag{6}
\]

另一方面，公式 (4) 蕴含

\[
c _ {\alpha , \beta , \gamma} = 0 \quad \text { if } | \alpha | + | \beta | <   | \gamma |\tag{7}
\]

并且展开式中总次数为 \(|\gamma|\) 的、\(z^\gamma\) 的项也正是 \((x_1 + y_1)^{\gamma_1}(x_2 + y_2)^{\gamma_2} \ldots (x_n + y_n)^{\gamma_n} = \sum_{\alpha + \beta = \gamma} ((\alpha, \beta)) x^\alpha y^\beta\) 的相应项（cf.~Differentiable and Analytic Manifolds, R, Notation and conventions）。因此：

\begin{enumerate}
\def\labelenumi{(\arabic{enumi})}
\setcounter{enumi}{7}
\tightlist
\item
\end{enumerate}

\[
c _ {\alpha , \beta , \gamma} = 0 \quad \text { if } \quad | \alpha | + | \beta | = | \gamma | \quad \text { but } \quad \alpha + \beta \neq \gamma\tag{8}
\]

\[
c _ {\alpha , \beta , \alpha + \beta} = ((\alpha , \beta)).\tag{9}
\]

乘法的结合律蕴含关系

\[
\sum_ {\alpha , \xi} c _ {\alpha \xi \eta} x ^ {\alpha} \left(\sum_ {\beta , \gamma} c _ {\beta \gamma \xi} y ^ {\beta} z ^ {\gamma}\right) = \sum_ {\xi , \gamma} c _ {\xi \gamma \eta} \left(\sum_ {\alpha , \beta} c _ {\alpha \beta \xi} x ^ {\alpha} y ^ {\beta}\right) z ^ {\gamma}
\]

对充分接近 0 的所有 \(x, y, z\)，从而

\[
\sum_ {\xi} c _ {\alpha \xi \eta} c _ {\beta \gamma \xi} = \sum_ {\xi} c _ {\xi \gamma \eta} c _ {\alpha \beta \xi} (\alpha , \beta , \gamma , \eta \text { in } \mathbf {N} ^ {n}).\tag{10}
\]

李群芽 G 有交换开子群芽，当且仅当 \(c_{\alpha\beta\gamma} = c_{\beta\alpha\gamma}\) 对所有 \(\alpha, \beta, \gamma\)（它们属于 \(\mathbf{N}^{n}\)）都成立。

